\section*{Rangkuman}
\addcontentsline{toc}{section}{Rangkuman}

Pada bab ini telah dibahas:
\begin{itemize}[leftmargin=*]
  \item peran lexer dan antarmuka pull model dengan parser;
  \item perbedaan token, lexeme, pattern, dan atribut token;
  \item definisi, operasi, prioritas, dan identitas ekspresi reguler;
  \item aturan maximal munch dan prioritas spesifikasi token;
  \item tabel ekspresi reguler untuk setiap kategori token;
  \item spesifikasi token Subset-C, header \file{token.h}, dan
        penanganan kesalahan leksikal.
\end{itemize}

Bab VI akan menghubungkan ekspresi reguler dengan mesin pengenal:
NFA, DFA, konversi NFA ke DFA, dan minimalisasi DFA. Teori ini menjadi
dasar implementasi lexer pada Bab VII.
